% 一纸成课 · 技术文档 PDF 排版样式（由 export-pdf.ps1 通过 --include-in-header 注入）
\usepackage{etoolbox}
\usepackage{fvextra}

% 代码块：字号略小、超长行自动折行，避免右侧溢出
\DefineVerbatimEnvironment{Highlighting}{Verbatim}{breaklines,commandchars=\\\{\},fontsize=\footnotesize}

% 表格：字号小一档、行距收紧，窄页面下仍能放下 4–5 列
\AtBeginEnvironment{longtable}{\footnotesize}
\renewcommand{\arraystretch}{0.92}
\setlength{\LTpre}{2pt}
\setlength{\LTpost}{4pt}

% 目录独占一页（只在 -Toc 时生效）
\apptocmd{\tableofcontents}{\clearpage}{}{}

% 图注与表注字号，图题用中文「图 N」
\usepackage{caption}
\captionsetup{font=small,labelfont=bf}
\renewcommand{\figurename}{图}
\renewcommand{\tablename}{表}

% 段落间距略收紧（正文 10 pt + 1.35 倍行距已由 pandoc 变量设置）
\setlength{\parskip}{2pt}
\setlength{\parindent}{0pt}

% 正文里少量符号（勾/叉/箭头/不等号）在主字体中缺字形，统一映射到数学符号，
% 只影响排版，不改 Markdown 源文件内容。
\usepackage{newunicodechar}
\newunicodechar{✅}{\ensuremath{\checkmark}}
\newunicodechar{✓}{\ensuremath{\checkmark}}
\newunicodechar{❌}{\ensuremath{\times}}
\newunicodechar{✕}{\ensuremath{\times}}
\newunicodechar{↔}{\ensuremath{\leftrightarrow}}
\newunicodechar{≤}{\ensuremath{\leq}}
\newunicodechar{≥}{\ensuremath{\geq}}
\newunicodechar{−}{\ensuremath{-}}
\newunicodechar{◐}{部分}
